% Top-to-bottom thesis structure with parallel research branches in grayscale.
\begingroup
\linespread{1}\selectfont
\begin{tikzpicture}[
  font=\thesisfigurefont,
  box/.style={draw=black!55,rounded corners=2pt,align=center,inner sep=4pt,line width=.6pt},
  wide/.style={box,minimum width=13.1cm,fill=black!5,minimum height=.8cm},
  branch/.style={box,minimum width=3.9cm,minimum height=2.7cm,fill=white},
  flow/.style={-{Latex[length=2mm]},draw=black!65,line width=.65pt},
  node distance=5mm
]
\node[wide] (problem) {\textbf{共同问题：任务持续演化与训练信息受限}};
\node[wide,below=of problem,minimum height=1.05cm] (framework)
  {\textbf{第三章：MedCL 评测平台设计与分割基准}\\数据与任务设定\quad 阶段矩阵\quad 多维评价指标};
\node[branch,below=1.05cm of framework] (fed)
  {\textbf{第五章：类别扩展}\\\textbf{联邦持续分类}\\[5pt]无原始样本回放\\子空间融合与梯度保护};
\node[branch,left=5mm of fed] (seg)
  {\textbf{第四章：域与类别增量}\\\textbf{弱监督持续分割}\\[5pt]有限回放\enspace ·\enspace 涂鸦监督\\监督增强与特征一致性};
\node[branch,right=5mm of fed] (reg)
  {\textbf{第六章：任务增量}\\\textbf{持续医学图像配准}\\[5pt]有限图像对回放\\元学习与锐度感知优化};
\node[wide,below=6mm of fed,minimum height=1.05cm,fill=black!8] (evaluation)
  {\textbf{共同评价：任务性能与持续学习能力}\\阶段矩阵\quad BWTR\quad RMA\quad MPE\quad DRR};
\node[wide,below=of evaluation] (future)
  {\textbf{第七章：总结与展望}\qquad 多模态持续学习};
\draw[flow] (problem) -- (framework);
\draw[flow] (framework.south -| seg.north) -- node[right,align=left] {数据与评价\\直接承接} (seg.north);
\draw[flow] (framework.south) -- node[right] {阶段矩阵与指标} (fed.north);
\draw[flow] (framework.south -| reg.north) -- (reg.north);
\draw[flow] (seg.south) -- (evaluation.north -| seg.south);
\draw[flow] (fed.south) -- (evaluation.north);
\draw[flow] (reg.south) -- (evaluation.north -| reg.south);
\draw[flow] (evaluation) -- (future);
\end{tikzpicture}
\endgroup
